% ============================================================
% Formulario, argomento 19: decisioni bayesiane e classificazione binaria (dispensa, capitolo 19)
% ============================================================
\argomento{Decisioni bayesiane e classificazione binaria}{Dispensa: capitolo 19}[Decisioni bayesiane]

\voce{regoladecisione}{Regola di decisione}{$M$ ipotesi $H_i$, ciascuna con la sua legge dei dati}
  {D : \vett{x} \in \mathcal{X}^n \;\longmapsto\; D(\vett{x}) \in \{1, 2, \ldots, M\}}
  {\variante{regioni di decisione, una partizione dei dati}{\Omega_i = \{\vett{x} : D(\vett{x}) = i\}}
   \variante{legge dei dati sotto $H_i$ e probabilità a priori}{p_{\vett{X}}(\vett{x} \mid H_i), \quad P(H_i), \quad
     i = 1, \ldots, M}}

\voce{rischio}{Rischio medio bayesiano}{$C_{ij}$: costo di decidere $i$ quando è vera $H_j$}
  {R = \sum_{i=1}^{M} \sum_{j=1}^{M} C_{ij}\, P\{D(\vett{X}) = i,\ H = H_j\}}
  {\variante{costi uniformi, $C_{ii} = 0$ e $C_{ij} = 1$ per $i \neq j$}{R = P(e)}
   \variante{la regola ottima}{\text{minimizza } R}}

\voce{perrore}{Probabilità di errore con due ipotesi}{si sbaglia decidendo per l'ipotesi falsa}
  {P(e) = P\{\vett{X} \in \Omega_1,\ H_2\} + P\{\vett{X} \in \Omega_2,\ H_1\}}
  {\variante{con le probabilità di errore condizionate}{P(e) = P(H_1)\, P(e \mid H_1) + P(H_2)\, P(e \mid H_2)}
   \variante{test con soglia, errore sotto $H_1$}{P(e \mid H_1) = P\{L(\vett{X}) < \eta \mid H_1\}}
   \variante{errore sotto $H_2$}{P(e \mid H_2) = P\{L(\vett{X}) > \eta \mid H_2\}}}

\voce{regolamap}{Regola MAP (minima probabilità di errore)}{si decide $H_1$ se il rapporto di verosimiglianza
  $L$ supera la soglia $\eta$}
  {L(\vett{x}) = \frac{p_{\vett{X}}(\vett{x} \mid H_1)}{p_{\vett{X}}(\vett{x} \mid H_2)} \decidi{H_1}{H_2}
   \frac{P(H_2)}{P(H_1)} = \eta}
  {\variante{senza rapporto}{p_{\vett{X}}(\vett{x} \mid H_1)\, P(H_1) \decidi{H_1}{H_2} p_{\vett{X}}(\vett{x} \mid
     H_2)\, P(H_2)}
   \variante{con Bayes: la massima probabilità a posteriori}{P\{H_1 \mid \vett{X} = \vett{x}\} \decidi{H_1}{H_2}
     P\{H_2 \mid \vett{X} = \vett{x}\}}
   \variante{dati continui}{L(\vett{x}) = \frac{f_{\vett{X} \mid H_1}(\vett{x} \mid H_1)}{f_{\vett{X} \mid
     H_2}(\vett{x} \mid H_2)} \decidi{H_1}{H_2} \frac{P(H_2)}{P(H_1)}}
   \variante{il simbolo $\decidi{H_1}{H_2}$ si legge: con $>$ si decide $H_1$, con $<$ si decide $H_2$}{}}

\voce{regolaml}{Regola ML (massima verosimiglianza)}{ipotesi equiprobabili: la MAP con $\eta = 1$}
  {L(\vett{x}) \decidi{H_1}{H_2} 1 \qquad\text{cioè}\qquad p_{\vett{X}}(\vett{x} \mid H_1) \decidi{H_1}{H_2}
   p_{\vett{X}}(\vett{x} \mid H_2)}
  {}

\voce{testhamming}{Test sul peso di Hamming}{due sorgenti di bit i.i.d., $p_1 > p_2$, ipotesi equiprobabili}
  {w_H(\vett{x}) \decidi{H_1}{H_2} \eta_1 = n\, \frac{\ln \dfrac{1 - p_2}{1 - p_1}}{\ln \dfrac{p_1 (1 - p_2)}{p_2 (1
   - p_1)}}}
  {\variante{verosimiglianza, $w_H$ numero di uni}{p_{\vett{X}}(\vett{x} \mid H_i) = p_i^{w_H(\vett{x})}\, (1 -
     p_i)^{n - w_H(\vett{x})}}
   \variante{errore sotto $H_1$, $\eta_1$ non intero}{P(e \mid H_1) = \sum_{i=0}^{\lfloor \eta_1 \rfloor}
     \binom{n}{i} p_1^i (1 - p_1)^{n - i}}
   \variante{errore sotto $H_2$}{P(e \mid H_2) = \sum_{i = \lfloor \eta_1 \rfloor + 1}^{n} \binom{n}{i} p_2^i (1 -
     p_2)^{n - i}}
   \variante{statistica sufficiente: basta il numero di uni $w_H(\vett{x})$}{}}

\voce{testmedia}{Test sulla media gaussiana}{$n$ dati i.i.d. $\mathcal{N}(\mu_i, \sigma^2)$, $\mu_1 > \mu_2$,
  ipotesi equiprobabili}
  {\frac{1}{n} \sum_{k=1}^{n} x_k \decidi{H_1}{H_2} \frac{\mu_1 + \mu_2}{2} \qquad P(e) = Q\left( \sqrt{n}\,
   \frac{\mu_1 - \mu_2}{2\sigma} \right) \xrightarrow[n \to \infty]{} 0}
  {\variante{densità dei dati sotto $H_i$}{f_{\vett{X} \mid H_i}(\vett{x} \mid H_i) = \prod_{k=1}^{n}
     \frac{1}{\sqrt{2\pi\sigma^2}}\, e^{-\frac{(x_k - \mu_i)^2}{2\sigma^2}}}
   \variante{$n = 1$, $\mu_{1,2} = \pm A$: la trasmissione binaria}{\text{soglia } 0, \quad P(e) = Q\left(
     \frac{A}{\sigma} \right)}
   \variante{statistica sufficiente: basta la somma dei dati}{}}
